\section{术语消化：本期关键词索引}

\begin{longtable}{p{0.24\textwidth}p{0.32\textwidth}p{0.35\textwidth}}
\toprule
术语 & 一句话解释 & 在本期中的作用 \\
\midrule
World Model & 根据状态和动作预测未来状态的模型 & 访谈的中心技术命题。 \\
Predictive Brain & 面向物理世界的预测性大脑 & AMI Labs 想构建的上游智能层。 \\
State & 对环境的任务相关压缩表示 & 连接表征学习与 planning。 \\
Action & 智能体对环境施加的动作或干预 & 让预测进入控制和决策。 \\
Representation Learning & 学习可复用的抽象表征 & 从 vision 通往 world model 的桥。 \\
Vision as Perspective & 把视觉理解成智能的基础视角，而非任务集合 & 解释 CV 在 LLM 时代的意义。 \\
LLM-pilled & 被 LLM 叙事过度组织资源和想象 & 对硅谷价值链的批评。 \\
VLA & Vision-Language-Action 模型 & 机器人路线的重要但不充分形态。 \\
World Model Pre-training & 面向视频/传感器/物理世界的基础预训练 & 预训练下半场的核心开放问题。 \\
Research Taste & 判断问题、证据、时机和审美的能力 & 解释个人和组织选择。 \\
JEPA & Joint Embedding Predictive Architecture & LeCun 路线中“在抽象表征空间预测”的代表思想。 \\
反向 OpenAI & 从世界和伙伴网络共建模型，而非只下载互联网 & AMI 的组织和商业隐喻。 \\
\bottomrule
\end{longtable}

\subsection{本章小结}

这些关键词都指向同一个问题：下一阶段 AI 是否能从“会处理人类写下来的世界”走向“能在真实世界中预测、行动和学习”。这需要模型、数据、组织和哲学边界同时更新。

